% Generated draft: requires mathematical and source review.
\documentclass[11pt]{article}
\usepackage[margin=25mm]{geometry}
\usepackage{fontspec}
\setmainfont{Latin Modern Roman}
\usepackage{amsmath,amssymb,mathtools,unicode-math}
\setmathfont{Latin Modern Math}
\usepackage{enumitem,needspace,xurl,hyperref}
\setlength{\parindent}{0pt}
\setlength{\parskip}{5pt}
\setlength{\emergencystretch}{3em}
\newcommand{\N}{\mathbb{N}}
\newcommand{\Z}{\mathbb{Z}}
\newcommand{\Q}{\mathbb{Q}}
\newcommand{\R}{\mathbb{R}}
\newcommand{\C}{\mathbb{C}}
\newcommand{\F}{\mathbb{F}}
\newcommand{\A}{\mathbb{A}}
\newcommand{\PP}{\mathbb{P}}
\newcommand{\E}{\mathbb{E}}
\newcommand{\eps}{\varepsilon}
\newcommand{\dd}{\,\mathrm{d}}
\begin{document}
\section*{UnsolvedMath ID: 11100029}
\textbf{Source problem code:} catalogue entry
\subsection{Definitions and mathematical statement}
The source asks to carry out Stolz's proposed programme for metrics of positive Ricci curvature on the class written ``\(MO8\) manifolds,'' with an expected obstruction in elliptic cohomology rather than \(K\)-theory. A Riemannian metric has positive Ricci curvature when \(\operatorname{Ric}(v,v)>0\) for every nonzero tangent vector. This is the Ricci condition in the source, not merely positive scalar curvature.

The manifold notation appears to have lost a connected-cover decoration. If it means \(MO\langle8\rangle\)-manifolds, the standard interpretation is string manifolds: their stable tangent classifying map lifts to the appropriate connected cover of \(BO\), equivalently spin data with a trivialization of the fractional class \(p_1/2\) in the usual convention. That is a tangential-structure condition in general dimension, not a claim that all manifolds have dimension eight. The exact source notation and tangential/normal convention must be checked.

A concrete obstruction programme would select an elliptic theory \(E\) and a compatible bordism orientation/genus, producing an invariant \(\alpha_E(M)\in\pi_{\dim M}E\), and establish a specified necessary vanishing condition for positive Ricci curvature. Whether a converse, surgery classification or complete obstruction is intended depends on the omitted Stolz plan; no necessary-and-sufficient theorem is inserted. Nonvanishing as an obstruction is also different from a numerical ``positivity'' property of a Witten genus.

The summary reports Witten-genus/tmf-orientation tools as partial progress but no established full elliptic obstruction to positive Ricci curvature, and refers to prior item \texttt{AMR-110-0028}. Its displayed \(M\operatorname{Spin}\to tmf\) orientation must not be adopted without the required extra structure or a different elliptic target: the usual full-tmf Witten-genus orientation is a string-structure construction. Its references to Wu/Gromov and Chodosh--Li give no exact Ricci-curvature theorem. The original ``hardly anything is known, I think'' is a tentative historical assessment, not a current verified absence result.
\subsection{Short English statement}
Develop the expected elliptic-cohomology obstruction programme for positive Ricci curvature on the intended connected-cover/string-type manifolds. The exact manifold notation, elliptic target and scope of Stolz's plan are missing, and scalar-curvature results cannot simply replace Ricci statements.
\subsection{Sources}
\begin{itemize}
\item Local UnsolvedMath registry, permanent ID \(11100029\), exact source code \texttt{catalogue entry}.
\item Mark Hovey, \emph{Hovey's Algebraic Topology Problem List}, Applications, item 5, surviving mirror: \url{https://www-users.cse.umn.edu/~tlawson/hovey/}. Registry access 2026-07-29; html-ordered-list extraction.
\item Stolz's plan is invoked without a full citation. The summary refers to \texttt{AMR-110-0028}, Witten genus/tmf, and unspecified Wu/Gromov and Chodosh--Li work, without precise bibliography or an exact theorem establishing the proposed obstruction.
\end{itemize}
\end{document}
